\title{LongRoPE: Extending LLM Context Window Beyond 2 Million Tokens}

\begin{abstract}
Having a large context window is highly advantageous in large language models (LLMs). Nevertheless, current models typically support context windows up to approximately 128k tokens because extending beyond this is hindered by expensive fine-tuning requirements, limited availability of lengthy training data, and instability caused by unfamiliar token positions.
\end{abstract}

In this work, we present LongRoPE, which, for the first time, successfully extends the context window of pre-trained LLMs up to a remarkable \textbf{2048k} tokens. This is accomplished using as few as 1k fine-tuning iterations confined to training sequences of length up to 256k, all while preserving the performance levels seen in the original short context window. Our approach relies on three principal contributions: (i) we detect and leverage two types of non-uniformities present in positional interpolation by employing a highly efficient search procedure, resulting in improved initialization for fine-tuning and facilitating up to an 8$\times$ increase in context length even when fine-tuning is not applied; (ii) we propose a stepwise extension mechanism, initially adapting the LLM to a 256k sequence length via fine-tuning, followed by a second-stage positional interpolation on this adapted LLM to reach a 2048k context window; (iii) we recalibrate LongRoPE on sequences of 8k length to ensure that performance on shorter contexts is not degraded. Comprehensive empirical evaluation of LLaMA2 and Mistral across a range of benchmark tasks corroborates the efficacy of our approach. Notably, LongRoPE allows models to preserve their original architecture, necessitating only minimal changes to the positional embeddings, and most established optimizations remain compatible. The implementation will be released at \texttt{https://github.com/microsoft/LongRoPE}
\end{abstract}

\section{Introduction}

Although Large Language Models (LLMs) have achieved impressive outcomes across a range of tasks~\cite{gpt4,llama2}, they are frequently constrained by their restricted context window—LLaMA2, for instance, has a limit of 4096 tokens~\cite{llama2}. When input exceeds this window, the model's effectiveness deteriorates, since it encounters positional information outside its training regime. This limitation becomes particularly problematic in critical applications such as in-context learning with many examples~\cite{Huang2023FewerIM} and LLM-driven agents~\cite{park2023generative,madaan2023selfrefine}.

Recent studies demonstrate that by fine-tuning pre-trained LLMs with longer sequences, the context window can be increased to approximately 128k tokens~\cite{longlora,pi,yarn,zhang2024soaring,liu2023scaling}. However, expanding the context window further presents three primary challenges. Firstly, the introduction of new and untrained positional indices brings in many problematic values, resulting in out-of-distribution complications and causing fine-tuning processes to converge poorly~\cite{pi}. This issue becomes particularly severe when scaling up from 4k to over $1000k$ positions, as over 90\% of these indices are novel. Secondly, effective fine-tuning typically necessitates access to textual data matching the target sequence lengths, yet such ultra-long texts—particularly those beyond $1000k$—are scarce in current datasets. In addition, the training on these very lengthy sequences demands substantial computation, entailing significant GPU time and resources. Thirdly, as the context window is stretched to extreme lengths, the model's attention is increasingly diluted across a vast array of token positions, which undermines its ability to perform well on shorter original context windows~\cite{pi}.

A commonly adopted method to address the initial challenge is the interpolation of RoPE positional embeddings~\cite{rope,pi}, whereby new positional indices are rescaled to fit within the original pretrained range, as depicted in Fig.\ref{fig:overview}. Position Interpolation (PI)~\cite{pi} employs a linear interpolation of RoPE’s rotary angles, determined by the extension ratio. In contrast, NTK~\cite{ntk,dynamicntk} proposes non-uniform interpolation and extrapolation procedures that vary across RoPE dimensions. YaRN~\cite{yarn} further subdivides RoPE dimensions based on frequency groups and applies extrapolation, NTK-based interpolation, or linear interpolation to each respective category. Nevertheless, the positional embedding mechanism within the Transformer demonstrates \emph{complex, non-uniform information entropy}. Current interpolation strategies fail to exploit this nuanced non-uniformity, resulting in information degradation and thereby restricting the attainable context window length.

Section~\ref{sec:analysis} empirically demonstrates two main observations: \textbf{(1)} Successful positional interpolation needs to address two types of non-uniformity: differences across RoPE dimensions and across token positions. Reduced interpolation is preferable for dimensions with lower RoPE values and for tokens near the sequence beginning, but the best approach varies according to the intended sequence length extension. 
\textbf{(2)} Factoring in these non-uniformities within positional interpolation enables preservation of information present in the original RoPE, especially within the most important dimensions and for early token positions. This reduces information loss introduced by the interpolation process, yielding improved initial conditions for subsequent fine-tuning. Furthermore, it permits a length extension of up to 8$\times$ in cases where further fine-tuning is not performed.

Building on these observations, we present \textbf{LongRoPE}, a robust approach that enlarges the context window of LLMs to exceed 2 \emph{million} tokens. The core of LongRoPE comprises three principal contributions. Firstly, {LongRoPE} leverages the inherent multidimensional variability in positional interpolation. Specifically, it determines suitable rescale factors for the rotation angles in each RoPE dimension by referencing token positions. Given that the parameter space for identifying these rescale factors increases exponentially with the desired extension ratio, {LongRoPE} employs an evolutionary search algorithm enhanced by two optimization strategies to substantially improve search efficiency. An illustrative example of the optimized, rescaled RoPE can be seen in Fig.~\ref{fig:overview}.

Subsequently, {LongRoPE} implements a resourceful and incremental approach to expand the context window to 2048k, circumventing the necessity for extensive fine-tuning on very long sequences, which are rarely available. The process initiates by identifying a suitable RoPE rescale for a 256k context length within the pre-trained LLM, followed by fine-tuning at this length. Afterward, owing to the non-uniform positional interpolation that enables an 8$\times$ increase without additional fine-tuning, a further search for optimal RoPE rescale factors is performed on the newly fine-tuned and length-extended LLM. Through this methodology, the 2048k context window is attained for both LLaMA2 and Mistral~\cite{mistral}.

Lastly, in order to address the potential drop in performance for the original (short) context window, {LongRoPE} further refines the RoPE rescale factors in the extended LLM. Following the approach used when increasing the context window from 256k to 2048k, we also apply scaling down to 4k and 8k windows on the LLM that was fine-tuned for 256k. This employs our search algorithm to minimize positional interpolation. At inference time, if the sequence length falls below 8k, RoPE is updated using the optimal rescale factors identified through our search.

Our comprehensive testing on multiple LLMs and diverse tasks involving extended context lengths validates the efficacy of our approach. The results indicate that LongRoPE consistently preserves low perplexity when evaluated from 4k up to 2048k sequence length, secures more than 90\% accuracy in passkey retrieval, and achieves performance similar to standard benchmarks created for a 4096-token context window. LongRoPE is compatible with all RoPE-embedding-based LLMs. The code and LongRoPE-2048k models are slated for release.

\section{Non-uniformity in Positional Interpolation}
\label{sec:analysis}

\subsection{Preliminary}

Transformer models require explicit positional information, often in the form of position embedding, to represent the order of input tokens. Our work focuses on the RoPE~\cite{rope} position embedding, which is widely used in recent  LLMs. For a token at position index $n$, its corresponding RoPE encoding can be simplified as follows:

\begin{equation}
		\fontsize{7.8}{7.8} \selectfont
	\label{eq:rope}
	\left[ cos(n\theta_0), sin(n\theta_0), cos(n\theta_1), \cdot\cdot\cdot, cos(n\theta_{d/2-1}), sin(n\theta_{d/2-1}) \right]   
\end{equation}

Here, $d$ denotes the embedding dimension, $n\theta_i$ refers to the rotary angle associated with the token positioned at $n$, and $\theta_i=\theta^{-2i/d}$ encodes the frequency of rotation. In RoPE, the standard base value assigned to $\theta$ is 10000.

\textbf{\emph{Context window extension ratio $s$} and positional interpolation}. The parameter $s$ is specified as the proportion of the new context length $L'$ compared to the initial context length $L$: $s=\frac{L'}{L}$.

To extend the context window from $L$ to $L'$, current positional interpolation methods suggest downscaling rotation frequencies $\theta_i$ based on extension ratio $s$. Let $\beta=\theta^{2/d}$, and $\lambda$ denote the actual rescale factor related to $s$, we   unify these positional interpolation methods as follows:

\begin{equation}	
	\fontsize{7.5}{7.5} \selectfont
	\label{eq:unified_rope}
	\begin{aligned}
		\left[cos\left(\frac{n}{\lambda(\beta)^0}\right), sin \left(\frac{n}{\lambda(\beta)^0}\right), cos\left(\frac{n}{\lambda(\beta)^1}\right), \cdot\cdot\cdot,  sin \left(\frac{n}{\lambda(\beta)^{d/2-1}}\right) \right]   
	\end{aligned}
\end{equation}

\textbf{Linear positional interpolation (PI)}. The PI method~\cite{pi} introduces a linear interpolation of positional indices inside the bounds of the pre-trained sequence length. Given a desired extension ratio $s$, all positional rotation angles are scaled down linearly by a factor $\lambda=s$ in each RoPE dimension. This approach, however, causes positional representations to become densely packed, making it challenging for the model to differentiate between tokens with nearby positions. As a result, PI typically exhibits suboptimal performance when used with large extension ratios.

\textbf{NTK-based interpolation and extrapolation}. The works of ~\cite{ntk,dynamicntk} approach RoPE through the lens of information encoding, leveraging Neural Tangent Kernel (NTK) theory~\cite{jacot2018neural,tancik2020fourier}. In order to address the issue of congested position representations within PI, they propose spreading the interpolation burden more evenly across the various RoPE dimensions. This is achieved by applying less scaling to the lower (higher frequency) dimensions and greater scaling to the higher (lower frequency) ones, thereby enabling both interpolation and extrapolation for positional encodings with $\lambda=s^{i}$. The dynamic NTK method~\cite{dynamicntk} introduces further refinement by modulating the extension ratio at each position according to the sequence’s current length. Unlike PI—which requires subsequent fine-tuning—NTK-based strategies are capable of enlarging the context window without the need for fine-tuning, although they are typically constrained to a maximum extension ratio of 4$\times$.

\textbf{YaRN}~\cite{yarn} offers a noteworthy advancement in positional interpolation by organizing RoPE dimensions into three distinct groups based on frequency. Each group is assigned a unique interpolation method: high-frequency dimensions are subjected to extrapolation ($\lambda$=1), low-frequency dimensions utilize linear interpolation (PI), and those with intermediate frequencies apply the NTK technique. The fundamental innovation of YaRN is its dimension grouping, which relies on empirically determined human selection. Consequently, this approach might not yield the best results when adapting to novel LLM architectures.

\subsection{Study on Non-uniform Positional Interpolation}

Drawing motivation from NTK and YaRN, we observe that their performance improvements stem largely from the introduction of non-linearity, notably through distinct frequency treatments across RoPE dimensions to facilitate tailored interpolation and extrapolation. Despite these advancements, the prevailing non-linear approaches are predominantly shaped by heuristics crafted by humans. This leads to the following inquiries: (1) Does existing positional interpolation represent the best possible approach? (2) Are there additional, uninvestigated forms of non-linearity?

In order to address these questions, we apply evolution search (refer to Sec.~\ref{sec:method}) to identify improved non-uniform positional interpolations specifically for LLaMA2-7B. The search process is directed by perplexity, utilizing 5 randomly selected samples from the PG19~\cite{pg19} validation dataset. Our empirical investigation leads us to several important conclusions.

\textbf{Finding 1}: \textit{RoPE dimensions exhibit substantial non-uniformities, which are not effectively handled by current positional interpolation methods.}

We optimize $\lambda$ values for each RoPE dimension as specified in Eq.~\ref{eq:unified_rope}. Table~\ref{tbl:exp1} presents a comparison of the perplexity results for LLaMA2-7B across several methods on the PG19 and Proof-pile~\cite{proof-pile} test sets, all without additional fine-tuning. Our optimized approach yields substantial perplexity reductions, highlighting that the widely used linear (PI) and non-uniform interpolation strategies (such as Dynamic-NTK and YaRN) do not reach optimal performance. In particular, YaRN achieves lower results compared to both PI and NTK on the PG19 test set; this is likely because YaRN fails to achieve the intended context window size in LLMs that have not been fine-tuned. For instance, when the context length is set to 8k, YaRN exhibits a sharp increase in perplexity after processing 7k tokens.

In our exploration, we observe that the rescaling coefficients $\lambda$ in Eq.~\ref{eq:unified_rope} are no longer uniform; unlike the constant scale $s$ used in PI, or the expressions computed in NTK and the group-based adjustments in YaRN, these factors vary across dimensions. This non-uniform scaling yields notable gains in LLaMA2's language modeling capabilities—measured by perplexity—at 8k and 16k context lengths, and does so without the need for additional fine-tuning. The underlying reason for this improvement is that such positional embeddings retain the structural characteristics of the original RoPE, especially in key dimensions, thereby making it easier for the language model to differentiate tokens that are close together in sequence.

\textbf{Finding 2}: \textit{RoPE for the initial tokens in the input sequence should be extrapolated with less interpolation.} 

We propose that the RoPE applied to the first $\hat{n}$ tokens of input sequences should involve less interpolation, as these tokens typically receive high attention scores and thus play a pivotal role within attention layers, a phenomenon also reported in Streaming LLM~\cite{streamingllm} and LM-Infinite~\cite{han2023lminfinite}. To investigate this, we expand the context window to 8k and 16k using PI and NTK, selectively omitting interpolation for the first $\hat n$ (0, 2, ..., 256) tokens. If $\hat n$=0, standard PI and NTK are employed. Table~\ref{tbl:tokenposition} presents two principal findings: \textbf{(1)} excluding position interpolation for the initial tokens significantly enhances performance; \textbf{(2)} the optimal value of $\hat n$ varies according to the desired context extension length.

\textbf{Finding 3}: \textit{Non-uniform positional interpolation effectively extends LLM context window in both fine-tuning and non-fine-tuning settings.} 

Although our experiments demonstrate that the searched non-uniform position interpolation method offers substantial improvements in extension performance at 8k and 16k context lengths without any further training, extending to even longer contexts necessitates fine-tuning. Therefore, we fine-tune LLaMA2-7B using our searched RoPE, targeting a context window of 64k tokens (refer to the Appendix for detailed configurations). According to Table~\ref{tbl:finetune}, our approach achieves notable performance gains over PI and YaRN, both prior to and following the fine-tuning of LLaMA2-7B. This advantage is attributed to the effective application of non-uniform positional interpolation, which reduces information loss and offers a superior initialization for the fine-tuning process.

\textbf{Summary}. Our research identifies two distinct non-uniformities: differential treatment of RoPE dimensions and diverse token positions. Harnessing these non-uniformities within positional interpolation leads to significant advancements in the ability of large language models to extend their context lengths.

\section{LongRoPE}

\label{sec:method}
Building upon these insights, this work proposes LongRoPE, which incorporates a search algorithm designed to leverage both types of non-uniformity. This approach enables the extension of the LLM context window to exceed 2 million tokens.

\subsection{Problem Formulation}

These two types of non-uniformity significantly expand the range of possible solutions and complicate the process of optimization. To mitigate these challenges, we recast the multidimensional non-uniform position interpolation optimization task into a search-based formulation.

For a LLM targeting a  context window size of $L'$ and lengthy input documents $\mathbf{X}$, where each $\mathbf{x}\in\mathbf{X}$ surpasses $L'$ in token length, we denote the original rotary angle of the $i^{th}$ dimension in RoPE embedding at token position $n$ as  $\frac{n}{\beta^i}$. The optimization problem is then formulated as follows:

\begin{equation}
\begin{gathered}
		\label{eq:problem}

\underset {\mathbf{x}\in\mathbf{X}; \, |\mathbf{x}|\ge L'}{ \operatorname {arg\,min} } \,  \mathcal{L}\, \left( \text{LLM} (\text{RoPE}, \mathbf{X})\right),	\text{where \;} 
\\
\scalebox{0.93}{
$\underset {\begin{subarray}{l} i=0,\cdot\cdot,\frac{d}{2}-1;\\ n\in[ 0,|\mathbf{x}|);\end{subarray}}{\text{RoPE}_(n)}= {\left[\cdot\cdot,\cos\left(\mathbb{I}(\hat{\lambda}_{i}, \hat{n})\times\frac{n}{\beta^i}\right),  \sin\left(\mathbb{I}(\hat{\lambda}_{i}, \hat{n})\times\frac{n}{\beta^i}\right),\cdot\cdot\right] }$}\\
	\text{where \;} \mathbb{I}(\hat{\lambda}_{i},\hat{n})=\begin{cases}
	1&\mbox{\;  $n<\hat{n}$}\\
{\frac{1}{\lambda}_{i}}&\mbox{\; $n\geq\hat{n}$}
\end{cases}
	\end{gathered}
\end{equation}

In this approach, a set of rescale factors $\mathbb{I}(\hat{\lambda}_{i},\hat{n})$ is introduced to address both types of non-uniformities. The parameters $\hat{\lambda_i}$ and $\hat{n}$ represent the non-uniformity found in RoPE dimensions and token position indices, respectively. More precisely, $\mathbb{I}(\hat{\lambda}_{i},\hat{n})$ serves to modify the rotation angle for the $i^{th}$ RoPE dimension, where $\hat{\lambda}_{i}$ acts as the rescaling coefficient and $\hat{n}$ functions as a threshold for token positions. For token positions less than $\hat{n}$, the original RoPE rotary angle $\frac{n}{\beta^i}$ remains unchanged since the rescale factor $\hat{\lambda}_{i}$ is not applied. Once token positions reach $n \geq \hat{n}$, the rescaling is enacted using the factor.

For a desired context window length $L'$, the goal is to determine the best rescaling factors ($\mathbb{I}(\hat{\lambda}_{0},\hat{n})$, $\mathbb{I}(\hat{\lambda}_{1},\hat{n})$, ... , $\mathbb{I}(\hat{\lambda}_{i},\hat{n})$, ...) corresponding to each of the $d$ RoPE dimensions. The intention is that, after adjusting the $\text{RoPE}$ encoding accordingly in the target $\text{LLM}$, the model attains the lowest possible loss $\mathcal{L}$ (i.e., perplexity) when predicting the next token for given inputs $\mathbf{X}$ of length $L'$.

\subsection{Searching the Non-uniform Position Interpolation}

To address the challenge outlined in Eq.~\ref{eq:problem}, we propose a straightforward but powerful approach that identifies the most suitable RoPE rescale factors. This method is designed to maximize the utilization of the complex, uneven characteristics present in multidimensional position embeddings.

\textbf{Search space}. We construct an expansive search space that incorporates both forms of non-uniformity. Table~\ref{tbl:searchspace} provides details of this search space. In particular, we permit the selection of distinct rescale factors for each RoPE embedding dimension. For clarity and ease in defining the search space, the optimization process targets $\lambda_i$ and $\hat{n}$, rather than directly searching for $\mathbb{I}(\hat{\lambda}_i, \hat{n})$; here, $\hat{\lambda}_i$ is given by $1/\lambda_i$. As outlined in Table~\ref{tbl:searchspace}, $\lambda_i$ ranges from 1.0 (representing simple extrapolation) up to $s \times 1.25$ (enabling broader interpolation compared to PI), with increments of 0.01, where $s$ denotes the ratio by which the target context window is extended.

The parameter $\hat{n}$ determines how many starting token positions are preserved using the standard RoPE embedding, meaning they are not affected by position interpolation. In practice, we conduct a search over $\hat{n}$ within the set \{0, 1, 2, 4, 8, 12, 16, 20, 24, 28, 32, 64, 128, 256\}. Setting $\hat{n}$ to 0 implies that every token position relies entirely on the rescale factors determined through search.

\textbf{Evolution-based search}. The search space outlined in Table~\ref{tbl:searchspace} encompasses a wide array of positional interpolation configurations, resulting in a substantial challenge for effective exploration. For instance, adopting a $s=4\times$ extension yields $400^{128/2}\times14=4\times10^{167}$ possible options. As the extension ratio increases, the search space grows at an exponential rate. To efficiently navigate this complex landscape, we employ evolution search~\cite{guo2020single}, complemented by two optimization strategies that significantly enhance search performance. The overall search process is depicted in Algorithm~\ref{alg:search}.

\textit{Enhanced initial population creation}. In place of using purely random initialization for a population containing $P$ rescale factors, we explicitly include the three RoPE rescale factors that align with PI, NTK, and YaRN as members of the starting population. The other $P$-3 individuals are produced by applying random mutations to these three rescale factors, each mutation occurring with probability $p$.

\textit{Monotonically non-decreasing constraint}. Once the initial population has been produced, we assess the LLM perplexity for every candidate. In detail, each individual’s designated RoPE rescale factors are incorporated into the target LLM to evaluate the perplexity for the input $\mathbf{X}$. Selection of parent candidates for further evolutionary steps is based on identifying the top-k individuals. Nonetheless, the expansive nature of the search space often results in random mutation and crossover operations traversing suboptimal regions, which can lead to superfluous perplexity calculations. This inefficiency is exacerbated for large $L'$, due to the significant computational overhead required for each perplexity evaluation.

To mitigate this issue, we enforce a monotonicity constraint, stipulating that the rescaled RoPE factors follow a non-decreasing sequence: $\lambda_i \leq \lambda_{i+1}$. Only those RoPE configurations adhering to this condition are utilized for evaluating perplexity in LLMs, which effectively narrows the scope of the search. More precisely, we mandate that $\lambda_i$ grows monotonically along the RoPE dimension indices ($i = 0, \ldots, 63$). This requirement is motivated by Neural Tangent Kernel (NTK) theory~\cite{jacot2018neural,tancik2020fourier,ntk}, which indicates that the lower-indexed dimensions—characterized by higher frequencies—necessitate reduced interpolation (thus, smaller $\lambda_i$), while dimensions associated with lower frequencies permit increased interpolation (therefore, larger $\lambda_i$ values).

\begin{algorithm}[t]
	\small
	\caption{The search algorithm}
	\textbf{Input:} target LLM, input samples $\mathbf{X}$, population size $P$, mutation size $N_1$, crossover size $N_2$, max iterations $\mathcal{T}$, mutation probability $p$\\

	\begin{algorithmic}[1]
		\label{alg:search}
		\STATE Top-k $\leftarrow$ $\phi$;	
		\STATE $\text{P}_0 \leftarrow$ \textit{Initialize\_population\_with\_optimization} ($P$, $\mathbf{X}$, $p$);
		\FOR{$i = 1$ to $\mathcal{T}$}
		\STATE \textit{Compute\_perplexity} (LLM, $\text{P}_{i-1}$, $\mathbf{X}$);
		\STATE Top-k $\leftarrow$ \textit{Update\_Topk} (Top-k, $\text{P}_{i-1}$);
		\STATE $\text{P}_{mutation} \leftarrow$ \textit{Mutation\_with\_mono\_constraint}(Top-k, $N_1$, $p$);
		\STATE $\text{P}_{crossover} \leftarrow$ \textit{Crossover\_with\_mono\_constraint}(Top-k, $N_2$);
		\STATE $\text{P}_i \leftarrow \text{P}_{mutation} \cup \text{P}_{crossover} \cup$ Top-k;
		\ENDFOR
		\STATE Return the member of Top-k exhibiting the minimum perplexity;
	\end{algorithmic}
\end{algorithm}

\textbf{8$\times$ extension achieved without fine-tuning}. Through the use of evolutionary search, we successfully pinpoint non-uniform RoPE rescale coefficients that maintain crucial dimensions and positional encoding, thereby reducing the information loss that typically arises from interpolation. Figure~\ref{fig:non-ft} illustrates how our technique enables LLaMA2’s context window to expand from 4k to 32k in the absence of fine-tuning. Conversely, current approaches like PI, as well as non-uniform NTK and YaRN, experience a sharp increase in perplexity after only a 2$\times$ extension.

\subsection{Extending LLM Context Window to 2048K}
\label{sec:extendto2048k}

\textbf{Progressively expanding to 2048k}. In this section, we present our approach for scaling the context window of pre-trained LLMs from the conventional 4k length up to beyond 2048k. Our non-uniform positional interpolation technique is capable of supporting up to an 8$\times$ increase in context length without necessitating any fine-tuning. When attempting significantly larger extensions, such as 512$\times$, fine-tuning becomes unavoidable. One possible strategy involves identifying suitable RoPE rescaling parameters specifically for the intended 2048k window and then applying fine-tuning. However, this approach is constrained by the high computational and resource demands associated with such extensive training. Additionally, our empirical observations suggest that effective fine-tuning at these high extension scales proves to be quite difficult (see Appendix).

Fortunately, LongRoPE demonstrates efficacy across both the base model and the extended version of the LLM after fine-tuning. As a result, we propose a streamlined, incremental approach that enables reaching the intended 2048k context window by employing only 1k fine-tuning steps, using a maximum training sequence length of 256k.

$\Diamond$ \textit{LongRoPE search applied for pre-trained LLM context expansion up to 256k}. Using LLaMA2 as a representative case, we perform search operations to reach desired context lengths of 128k and 256k. In this phase, the enlargement rates are 32$\times$ and 64$\times$ for the respective window sizes.

$\Diamond$ \textit{Fine-tuning to 256k}. Next, the pre-trained LLM undergoes fine-tuning aimed at extending the context window to 256k. The process begins with LLaMA2 being fine-tuned for 400 steps using RoPE rescaling factors adjusted for 128k. Once this phase is complete, the RoPE rescaling factors are updated to 256k on the obtained checkpoint, followed by a further 600 steps of fine-tuning. This sequential approach demonstrates greater efficiency compared to direct fine-tuning for 256k.

$\Diamond$ \textit{Expanding fine-tuned extended LLM to 2048k via LongRoPE search}. In the final stage, we apply an additional search process to the LLM that has already been fine-tuned for 256k context length. Through this method, the model's context window is dramatically increased to 2048k, all without requiring extra rounds of fine-tuning. This achieves a total extension factor of $512\times$.

\textbf{Mitigating performance loss in shorter context windows}.  Upon expanding the context window to a vast 2048k tokens, we observe a reduction in effectiveness within the model's initial context range. This phenomenon is documented for positional interpolation methods~\cite{pi}, which confine position embeddings for the original window to a restricted region of a higher-dimensional space. Such compression hampers the language model’s capability. In scenarios with a 512$\times$ increase, the positions in the original 4k context window are densely packed, exacerbating this issue.

To address this issue, we conduct an additional evolution search on the expanded LLM to fine-tune the RoPE rescale factors tailored for shorter context lengths, such as 4k and 8k. For these shorter contexts, we constrain the upper bound of the search space for $\lambda$, since less positional interpolation is necessary. At inference time, the LLM automatically modifies the relevant RoPE rescale factors based on the input.

\section{LongRoPE}

\label{sec:method}
Building upon the insights obtained, we propose LongRoPE. This approach incorporates an optimized search algorithm designed to leverage both identified non-uniformities, enabling the extension of the LLM context window to surpass 2 million tokens.

\subsection{Problem Formulation}

These two types of non-uniformity result in an expansive solution space and complicate the optimization process. To tackle this challenge, we recast the multidimensional non-uniform position interpolation optimization as a search problem.

For a LLM targeting a  context window size of $L'$ and lengthy input documents $\mathbf{X}$, where each $\mathbf{x}\in\mathbf{X}$ surpasses $L'$ in token length, we denote the original rotary angle of the $i^{th}$ dimension in RoPE embedding at token position $n$ as  $\frac{n}{\beta^i}$. The optimization problem is then formulated as follows:

\begin{equation}
\begin{gathered}
		\label{eq:problem}

\underset {\mathbf{x}\in\mathbf{X}; \, |\mathbf{x}|\ge L'}{ \operatorname {arg\,min} } \,  \mathcal{L}\, \left( \text{LLM} (\text{RoPE}, \mathbf{X})\right),\;\text{with} 
\\
\scalebox{0.93}{
$\underset {\begin{subarray}{l} i=0,\cdot\cdot,\frac{d}{2}-1;\\ n\in[0,|\mathbf{x}|);\end{subarray}}{\text{RoPE}_{(n)}} = {\left[\ldots, \cos\left(\mathbb{I}(\hat{\lambda}_{i}, \hat{n}) \frac{n}{\beta^i}\right), \sin\left(\mathbb{I}(\hat{\lambda}_{i}, \hat{n}) \frac{n}{\beta^i}\right), \ldots\right]}$}
\\
	\text{with} \;\; \mathbb{I}(\hat{\lambda}_{i},\hat{n})=\begin{cases}
	1 & \text{if}\;\; n < \hat{n} \\
	{\frac{1}{\lambda}_{i}} & \text{if}\;\; n \geq \hat{n}
	\end{cases}
	\end{gathered}
\end{equation}

In this approach, we define a collection of rescaling factors $\mathbb{I}(\hat{\lambda}_{i},\hat{n})$ designed to address two distinct types of non-uniformities. Here, $\hat{\lambda_i}$ represents the non-uniformity corresponding to specific RoPE dimensions, while $\hat{n}$ pertains to disparities across token positions. The function $\mathbb{I}(\hat{\lambda}_{i},\hat{n})$ is employed to adjust the rotation angle associated with the $i^{th}$ dimension of RoPE: the rescaling parameter $\hat\lambda_{i}$ and the position threshold $\hat{n}$ govern this transformation. For token positions less than $\hat{n}$, namely the initial $\hat{n}$-1 positions, $\hat\lambda_{i}$ does not modify the rotary angle, thus retaining the standard RoPE formula $\frac{n}{\beta^i}$. When token indices satisfy $n\ge\hat{n}$, the rescaling factor is incorporated.

With a target context window length set at $L'$, our goal is to determine the set of optimal rescale factors—represented as ($\mathbb{I}(\hat{\lambda}_{0},\hat{n})$, $\mathbb{I}(\hat{\lambda}_{1},\hat{n})$, ..., $\mathbb{I}(\hat{\lambda}_{i},\hat{n})$,...)—across each of the $d$ RoPE dimensions. By applying these rescaling coefficients, we seek to enable the target $\text{LLM}$ employing the modified $\text{RoPE}$ to attain the lowest possible next-token prediction loss, denoted by $\mathcal{L}$ (corresponding to perplexity), when processing input sequences $\mathbf{X}$ of token length $L'$.

\subsection{Searching the Non-uniform Position Interpolation}

In order to address the challenge outlined in Eq.~\ref{eq:problem}, we propose a straightforward but powerful approach that identifies the optimal RoPE rescale parameters, thereby maximizing the utilization of multidimensional variations within the position embedding.

\textbf{Search space}. We construct an extensive search space that encompasses both types of non-uniformity. The search space is outlined in Table~\ref{tbl:searchspace}. More concretely, we permit individualized optimization of the rescale factor for each dimension in the RoPE embedding. To streamline the design process, the search is performed over $\lambda_i$ and $\hat{n}$ rather than directly searching for $\mathbb{I}(\hat{\lambda}_{i},\hat{n})$, where  $\hat{\lambda}_{i}=1/\lambda_i$. As detailed in Table~\ref{tbl:searchspace}, $\lambda_i$ ranges from a minimum of 1.0 (representing straightforward extrapolation) up to a maximum of $s\times1.25$ (corresponding to broader interpolation than PI), using increments of 0.01. Here, $s$ denotes the intended expansion ratio of the context window.

The parameter $\hat{n}$ determines how many starting token positions preserve the original RoPE embedding, foregoing position interpolation. In practice, we permit $\hat{n}$ to be selected from the set \{0, 1, 2, 4, 8, 12, 16, 20, 24, 28, 32, 64, 128, 256\}. If $\hat{n}=0$, every token position employs the optimized rescale factors.

\textbf{Evolution-based search}. The search space outlined in Table~\ref{tbl:searchspace} encompasses a vast array of positional interpolation strategies, creating substantial obstacles for effective navigation. For instance, extending to $s=4\times$ results in $400^{128/2}\times14$=4$\times10^{167}$ possible options. The scope of the search increases exponentially with higher extension ratios. To mitigate these complexities, we leverage evolution search~\cite{guo2020single} alongside two optimization strategies designed to markedly enhance search speed. The comprehensive search process is presented in Algorithm~\ref{alg:search}.

\textit{Enhanced initialization of the population}. In place of a fully random initialization for the $P$ rescale factors, we explicitly incorporate the three RoPE rescale factors associated with PI, NTK, and YaRN as members of the initial population. The other $P$-3 members are produced by applying random mutations to these three rescale factors, with a mutation probability set to $p$.

\textit{Monotonically non-decreasing constraint}. Following the creation of the initial population, we evaluate the LLM perplexity for every candidate. This involves applying each candidate's RoPE rescale factors to the target LLM and determining the perplexity for the input $\mathbf{X}$. The best $k$ candidates are chosen as parents to proceed with evolution. Nevertheless, the immense size of the search space makes simple mutation and crossover likely to traverse suboptimal regions, resulting in redundant perplexity evaluations. Such inefficiency is amplified in scenarios where $L'$ is large, as each perplexity assessment requires considerable inference time.

To tackle this issue, we enforce a constraint ensuring the RoPE rescaling factors are monotonically non-decreasing: $\lambda_i\le \lambda_{i+1}$. Only those RoPE configurations adhering to this condition are utilized for perplexity measurements in LLMs, which drastically decreases computational expense during the search process. In particular, we stipulate that the sequence $\lambda_i$ should rise monotonically with the RoPE dimension index ($i=0,\ldots,63$). This requirement is grounded in NTK theory~\cite{jacot2018neural,tancik2020fourier,ntk}, which argues that dimensions with greater frequency (i.e., lower indices) necessitate minimal interpolation (implying a smaller $\lambda_i$), while dimensions exhibiting lower frequency (i.e., higher indices) are capable of greater interpolation (thus allowing for a larger $\lambda_i$).

\begin{algorithm}[t]
	\small
	\caption{The search algorithm}
	\textbf{Input:} target LLM, input samples $\mathbf{X}$,   population size $P$, mutation size $N_1$, crossover size $N_2$, max iterations $\mathcal{T}$, mutate probability $p$\\

	\begin{algorithmic}[1]
		\label{alg:search}
		\STATE Initialize Top-k as an empty set;	
		\STATE Set $\text{P}_0$ to \textit{Initialize\_population\_with\_optimization} ($P$, $\mathbf{X}$, $p$); 
		\FOR{$i$ from $1$ to $\mathcal{T}$}
		\STATE Execute \textit{Compute\_perplexity} (LLM, $\text{P}_{i-1}$, $\mathbf{X}$);
		\STATE Update Top-k using \textit{Update\_Topk} (Top-k, $\text{P}_{i-1}$);
		\STATE Obtain $\text{P}_{mutation}$ by calling \textit{Mutation\_with\_mono\_constraint} (Top-k, $N_1$, $p$); 
		\STATE Create $\text{P}_{crossover}$ via \textit{Crossover\_with\_mono\_constraint} (Top-k, $N_2$);
		\STATE Combine $\text{P}_{mutation}$, $\text{P}_{crossover}$, and Top-k to form $\text{P}_i$;
		\ENDFOR
		\STATE Output the member with minimal perplexity in Top-k;
	\end{algorithmic}
\end{algorithm}

\textbf{8$\times$ extension without fine-tuning}. By leveraging evolutionary search, we efficiently determine optimal non-uniform RoPE rescale factors that retain critical dimensions and positional information, thereby reducing the information loss typically associated with interpolation. As illustrated in Fig.\ref{fig:non-ft}, this approach enables LLaMA2's context window to be increased from 4k up to 32k tokens without requiring any fine-tuning. In comparison, alternative methods like PI, non-uniform NTK, and YaRN exhibit significant increases in perplexity when the context window is extended beyond 2$\times$.

\subsection{Extending LLM Context Window to 2048K}
\label{sec:extendto2048k}

\textbf{Progressive extension to 2048k}. In this section, we present our approach for expanding the context window of pre-trained LLMs beyond the standard 4k limit, reaching up to 2048k. As shown, utilizing our non-uniform positional interpolation strategy enables an 8$\times$ context extension with no need for fine-tuning. However, for much larger extensions, such as up to 512$\times$, fine-tuning becomes essential. A possible strategy involves identifying suitable RoPE rescaling factors at the target length of 2048k, followed by subsequent fine-tuning. Nevertheless, this approach encounters significant obstacles, primarily due to the extremely high computational costs associated with such training. Additionally, our observations indicate that achieving effective fine-tuning is particularly difficult when the extension ratio is so large (refer to Appendix).

Fortunately, LongRoPE demonstrates efficacy with both baseline and fine-tuned versions of the extended LLM. To this end, we propose an efficient and incremental approach that reaches the target 2048k context length utilizing only 1k fine-tuning steps and a maximum training sequence length of 256k.

$\Diamond$ \textit{LongRoPE-based expansion of pre-trained LLMs to 256k: search process}. Using LLaMA2 as a case study, we perform a search procedure to achieve context window sizes of 128k and 256k tokens. At these stages, the respective extension ratios are 32$\times$ and 64$\times$.

$\Diamond$ \textit{Fine-tuning to 256k}. In the subsequent phase, we adapt the previously trained LLM to extend its context window to 256k. The procedure begins with 400 fine-tuning steps on LLaMA2 utilizing the RoPE rescaled factors corresponding to 128k. After completing these steps, we switch the RoPE rescaled factors to 256k on the obtained checkpoint and proceed with another 600 fine-tuning steps. This approach is found to be more efficient compared to direct fine-tuning towards 256k.

$\Diamond$ \textit{Expanding fine-tuned extended LLM to 2048k using LongRoPE search}. In the last phase, we apply an additional search procedure to the fine-tuned LLM configured with a 256k context length. As a result, the model achieves a substantial increase in context window to 2048k, attained without additional fine-tuning. This corresponds to a total extension factor of $512\times$.

\textbf{Restoring performance in smaller context windows}. Upon increasing the context window to a very large 2048k length, we observe that the model's performance deteriorates inside its initial context window. This phenomenon has been previously documented with positional interpolation~\cite{pi}, which compresses position embeddings in the original context window into a limited subspace due to the expansion into higher dimensions. Such compression impairs the effectiveness of the language model. Specifically, when the extension ratio is 512$\times$, the positions contained within the primary 4k context window are extremely densely packed.

To address this issue, we conduct an additional evolution search on the expanded LLM in order to fine-tune the RoPE rescale parameters for shorter context lengths (such as 4k and 8k). The upper limit for the searched $\lambda$ is lowered, since positional interpolation is less needed with shorter contexts. In inference, the LLM automatically updates the relevant RoPE rescale parameters as required.

 \section{Experiments}

\label{sec:exp}
\subsection{Setup}
\textbf{Evaluation Tasks and Models}. LongRoPE is utilized with both LLaMA2-7B and Mistral-7B architectures. Assessment is conducted in three areas: (1) the perplexity levels exhibited by large language models when processing documents with extended context; (2) performance on the Passkey retrieval challenge, which evaluates how effectively a model can locate a straightforward passkey amidst a large volume of extraneous information; and (3) results on conventional LLM benchmark tests, employing a limited context window of 4096 tokens.

\textbf{Fine-tuning}. For LLaMA2, the model is fine-tuned using a learning rate of 2e-5 that linearly decays, employing a global batch size of 32. Training is conducted for 400 steps on the Redpajama~\cite{redpajama} dataset, which is divided into 128k-length segments and uses BOS and EOS tokens as boundaries. Building on this initial checkpoint, an additional 600 steps are performed to extend the context window to 256k. Training with a 128k context window leverages 8 A100 GPUs, utilizing the distributed training framework~\cite{cube}, while scaling up to a 256k window necessitates the use of 16 A100 GPUs. For the Mistral model, the learning rate remains constant at 1e-6, with a global batch size of 64. Both 128k and 256k context size models adopt the YaRN~\cite{yarn} protocol: they undergo 400 steps of training on Together Computer's Long-Data Collections~\cite{mistraldata} with a 16k sequence length, using 4 A100 GPUs throughout the process.

\textbf{Search}. For cases where the target window size does not exceed 256k, the parameters are set as follows: $P$=64, $N_1$=$N_2$=16, $p$=0.3, and $\mathcal{T}$=40. In every iteration, the top 32 candidates are retained for mutation and crossover procedures. Perplexity assessment relies on five randomly selected samples from the PG19 validation set, each required to meet or exceed the specified target context length. When the window size surpasses 512k, the quantities for population, mutation, and crossover are each reduced by half. In this case, perplexity is evaluated using three randomly chosen samples from the Pile-Books3~\cite{pile} validation data.

\textbf{Baselines}. To achieve context windows of 2048k tokens, we fine-tuned models initially trained with 128k and 256k window sizes. This results in LongRoPE-2048k (ft=128k) for LLaMA2 and LongRoPE-2048k (ft=256k) for Mistral. Our evaluation contrasts these four models with leading approaches for expanding context windows, focusing on publicly available LLMs further trained after positional interpolation using PI, NTK, and YaRN techniques. The comparison encompasses Together-32k~\cite{together}, Code LLaMA~\cite{codellama}, LongLoRA-full-FT-100k~\cite{longlora}, YaRN-LLaMA, and YaRN-Mistral~\cite{yarn}.

\subsection{Main Results}

\label{sec:mainresult}
\textbf{Language modeling on extended sequences up to 256k tokens}. To assess our approach, we benchmark against leading LLMs optimized for longer contexts within the 256k token evaluation window. For a comprehensive assessment of robustness, we utilize both the Proof-pile~\cite{pg19} and PG19~\cite{pile} test datasets. Perplexity is measured at multiple context window sizes using a sliding window approach of 256 tokens. In the case of PG19, the evaluation encompasses the entire test split, which comprises 100 documents. For Proof-pile, consistent with YaRN~\cite{yarn}, we randomly select 10 samples, ensuring each sample exceeds 128k tokens in length.

Tables~\ref{tbl:proof-pile} and~\ref{tbl:pg19} present a comparison of the perplexity metrics for LLaMA2 and Mistral models that have been extended using various interpolation strategies on the Proof-pile and PG19 datasets, respectively. From these results, we draw two principal conclusions: \textbf{(1)} the extended models consistently exhibit reduced perplexity scores as the evaluation length increases from 4k to 256k, demonstrating their improved capacity to utilize extended context information. \textbf{(2)} Notably, the LongRoPE-2048k variants achieve superior performance compared to leading baseline models on context lengths up to 256k, despite operating with a context window that is 16$\times$ longer than standard, a scenario typically associated with significant performance degradation at shorter contexts.

\textbf{Language modeling on sequences exceeding 2000k}. To assess model performance on ultra-long textual data, we utilize the Books3~\cite{pile} corpus. For the sake of efficient evaluation, we draw a random sample of 20 books, all with lengths surpassing 2048k, and implement a sliding window of 256k tokens.

According to Table~\ref{tbl:books3}, LongRoPE is able to scale the context window of both LLaMA2-7B and Mistral-7B up to 2048k tokens, while maintaining perplexity on par with or better than baseline models for shorter sequences ranging from 8k to 128k. The comparison between the 2048k variants of LLaMA2 and Mistral reveals distinct behavioral patterns. Mistral demonstrates superiority over baselines for short contexts, but its perplexity surpasses 7 when the input length exceeds 256k. As anticipated, LLaMA2 exhibits steadily decreasing perplexity as context length increases, aside from minor rises at 1024k and 2048k. Further, the LLaMA2 model shows improved performance for LongRoPE-2048k when fine-tuned at 256k, as opposed to 128k, owing to a smaller secondary extension ratio (8$\times$ versus 16$\times$). By contrast, Mistral achieves higher performance with a fine-tuning window of 128k. This difference is largely attributable to the fact that, for Mistral's fine-tuning at 128k and 256k, the YaRN protocol uses a 16k training context, which influences Mistral's capacity to extend its window length in post-fine-tuning evaluations.

\textbf{Passkey retrieval}. In this section, we examine the practical context window capacity for generation tasks. We adopt a synthetic evaluation method called passkey retrieval, as described by~\cite{passkey}. This approach involves prompting the model to locate a randomly chosen passkey (a five-digit code) embedded within an extensive document. The specific prompt structure is outlined in the appendix. For this experiment, we conduct 10 runs of the passkey retrieval task, assigning the passkey to a randomly selected position uniformly distributed throughout the entire evaluation context.

Fig.~\ref{fig:passkey} illustrates how retrieval accuracy compares across different baselines. The accuracy of previous models declines sharply to zero once the sequence length surpasses 128k tokens. On the other hand, LongRoPE-LLaMA2-2048k (ft=256k) demonstrates robust performance, successfully retaining high retrieval accuracy ($\ge$90\%) over a wide range of sequence lengths, from 4k up to 2048k tokens—even under the highly demanding task of extracting a passkey from millions of tokens. Similarly, LongRoPE-Mistral-2048k (ft=128k) sustains perfect accuracy (100\%) through 1800k tokens, only decreasing to 60\% at 2048k. This result is consistent with observations in Table~\ref{tbl:books3}, where perplexity slightly rises at the 2048k mark.

\textbf{Evaluation on standard benchmarks within the original context length}. We assess the performance of LongRoPE-2048k models in their native context window utilizing the Hugging Face Open LLM Leaderboard~\cite{huggingfaceboard}, under both zero-shot and few-shot paradigms. Specifically, we employ 25-shot ARC-Challenge~\cite{allenai:arc}, 10-shot HellaSwag~\cite{zellers2019hellaswag}, 5-shot MMLU~\cite{mmlu}, and 0-shot TruthfulQA~\cite{lin2021truthfulqa}.

According to Table~\ref{tbl:commonsense}, our models perform similarly to the benchmarks intended for limited context windows and surpass the original Mistral model on TruthfulQA with an improvement of +0.5\%. Although LongRoPE-LLaMA2-2048k—fine-tuned at 256k—experiences somewhat greater performance decline, its results for most tasks are still considered satisfactory.

\subsection{Ablation Results}

\textbf{Evaluating the second positional interpolation's impact}. Within the framework of our progressive extension approach, the search algorithm facilitates a second application of non-uniform positional interpolation to fine-tuned extended LLMs. To assess its benefits, we experimented with the fine-tuned LLaMA2-256k model, expanding its context length to 512k, 1024k, and 2048k through both PI and YaRN methods. The data in Table~\ref{tbl:secondsearch} demonstrate that our non-uniform positional interpolation consistently maintains perplexity across different extension ratios. By comparison, perplexity under both PI and YaRN rises markedly as the context window is enlarged.

\textbf{Assessing recovery efficiency at reduced context lengths.} In order to address performance degradation when operating at shorter contexts, we modified the RoPE scaling parameters for LongRoPE-2048k using our optimization routine. In particular, we constrained the search process by lowering the maximum scale factors, which minimizes interpolation for brief contexts such as 4k and 8k tokens. Table~\ref{tbl:shortperformance} presents the perplexity outcomes for LongRoPE-LLaMA2-2048k evaluated on Proof-pile at 4k and 8k tokens, together with the mean accuracy on large language model benchmarks. The evidence underscores notable enhancements in performance at limited context lengths.

\textbf{Examination of the two types of non-uniformities}. In order to determine the individual effect of each non-uniformity on model performance, we conducted an ablation study. This involved two separate experimental designs: (i) lengthening LLaMA2-7B to 16k and 32k tokens utilizing three distinct strategies—PI, optimization of RoPE dimension exclusively, and optimization of both non-uniformities; (ii) scaling up a 256k-length LLaMA2 (after fine-tuning) to 2048k with a consistent approach. Perplexity values were assessed in the absence of additional fine-tuning. As illustrated in Table~\ref{tbl:searchdimension}, adapting non-uniformity within the RoPE dimension considerably lowers perplexity compared with PI's conventional linear interpolation. Incorporating non-uniformity based on token position yields clear gains for shorter sequences (16k and 32k tokens); however, such improvement is not evident at 2048k tokens—potentially attributable to the excessive sequence length. Simply retaining the original tokens, omitting interpolation, fails to provide benefit under these circumstances. Further investigation into this aspect is proposed as future work.

\section{Related Works}

Besides position interpolation techniques, this section reviews alternative strategies described in related literature.

\textbf{Retrieval-based} approaches rely on an external memory system to store extended past contexts and employ retrieval mechanisms to access relevant documents during inference~\cite{longllama,wang2023augmenting,borgeaud2022improving}. These methods frequently demand substantial, explicit adjustments to the underlying LLM architectures. In comparison, our method remains notably lightweight, introducing only minimal alterations to the positional embeddings. Furthermore, our approach enables support for a broader range of long-context applications beyond simple retrieval, such as extended document summarization and few-shot learning scenarios.

\textbf{Attention-based context window extensions}. In addition to interpolating positional embeddings, certain studies have expanded the effective input context of LLMs—without increasing the context window size—through modifications to the attention mechanism~\cite{han2023lminfinite, streamingllm,parallelwindow}. Central to this strategy is the reduction of the computational burden from newly introduced positions via innovative attention masks. Notably, these techniques complement positional interpolation approaches.

\textbf{Fine-tuning based approaches} investigate strategies for efficiently adapting pre-trained LLMs to handle extended context lengths through modifications to position embeddings. For instance, methods such as Code LLaMA~\cite{codellama}, LLaMA2 Long~\cite{xiong2023effective}, and ScaledRoPE~\cite{liu2023scaling} set the RoPE base value to be exceptionally large and conduct fine-tuning for the intended sequence length. In contrast, our approach enables compatibility with a range of target lengths and demonstrates the ability to scale beyond lengths of 2M. Additionally, since fine-tuning at very long context lengths (i.e., greater than 128k) is computationally intensive and requires significant GPU capacity, techniques like LongLoRA~\cite{longlora} and PoSE~\cite{zhu2023pose} have recently been introduced to address these resource challenges. Notably, our technique is complementary to these resource-efficient fine-tuning methods.

\section{Conclusion}

This paper introduces LongRoPE, a technique that substantially increases the context window of large language models (LLMs) to an extraordinary 2048k tokens, all the while preserving their performance within the original, shorter context range. Our approach leverages two distinct non-uniformities in RoPE positional embeddings, identified through an efficient evolutionary search process. This approach yields dual advantages: it facilitates robust initialization for subsequent fine-tuning and allows for an 8$\times$ expansion of the context window without the need for fine-tuning. Further, we present a step-wise extension strategy that employs LLMs fine-tuned on a 256k token context to achieve a 2048k context length without additional fine-tuning. Comprehensive experiments demonstrate the efficacy of LongRoPE. We anticipate that LongRoPE-2048k models will unlock novel long-context applications and stimulate further advancements in the field.

\section*{Broader Impacts} The objective of this study is to further the development of Machine Learning as a research area. While our contributions may yield various societal outcomes, we do not identify any particular impact that warrants special attention in this context.

	\balance

\end{document}